%% Ecological Modelling (Elsevier)
%% Build PDF: pdflatex -> bibtex -> pdflatex -> pdflatex
%% Build DOCX: ./make_docx.sh
\documentclass[preprint,12pt,authoryear]{elsarticle}

\usepackage{amsmath,amssymb}
\usepackage{graphicx}
\usepackage{booktabs}
\usepackage{multirow}
\usepackage[T1]{fontenc}
\usepackage[utf8]{inputenc}
\usepackage[hidelinks,bookmarks=false,pdfpagemode=UseNone]{hyperref}
\hypersetup{colorlinks=false,pdfborder={0 0 0}}
\usepackage{orcidlink}

%% Typography. microtype removes almost every overfull line by protrusion and
%% font expansion alone, without altering a single word. \file wraps a script
%% name in \path, which may break after its underscores, and \url does the same
%% for a bare address. Those two, plus the tables below, were the whole problem.
\usepackage{microtype}
\usepackage{url}
\urlstyle{same}
\def\UrlBreaks{\do\_\do\/\do\.\do\-}
\newcommand{\file}[1]{\path{#1}}

%% A little extra stretch lets TeX settle the few remaining prose lines instead
%% of pushing them into the margin, and slightly tighter column gutters keep the
%% widest tables inside the text block.
\setlength{\emergencystretch}{3em}
\setlength{\tabcolsep}{4.5pt}

\journal{Ecological Modelling}

\newcommand{\Mgha}{Mg\,ha\textsuperscript{-1}}
\newcommand{\eqnum}[1]{\qquad\qquad\mathrm{(#1)}}

\begin{document}

\begin{frontmatter}

\title{Prediction-interval coverage in spatial predictive mapping: design-based
assessment and correction under clustered calibration}

\author[uf]{Hugo Rodrigues\,\orcidlink{0000-0002-8070-8126}\corref{cor1}}
%% elsarticle echoes the whole \author argument in the e-mail footnote, which
%% would drag the ORCID mark down with it. Hand it just the name.
\makeatletter\@eadauthor={Hugo Rodrigues}\makeatother
\ead{rodrigues.machado.hugo@gmail.com}

\author[uf]{Golmar Golmohammadi\,\orcidlink{0000-0001-5532-3892}}

\address[uf]{University of Florida, Range Cattle Research and Education
Center, 3401 Experiment Station, Ona, FL 33865, USA}

\cortext[cor1]{Corresponding author.}

\begin{abstract}
Spatial predictive maps of environmental variables increasingly carry a
prediction interval alongside each value, and those intervals now inform carbon
accounting. Map accuracy must be estimated from a probability sample by
design-based inference, yet interval coverage is rarely confronted with the
sampling design at all. We show that coverage is a
population parameter of the same kind as the mean squared error, so the
Horvitz-Thompson argument that makes accuracy design-unbiased applies to it
unchanged, and is here exactly rather than approximately unbiased. On the
Amazonian above-ground biomass benchmark of Wadoux et al. (2021), CV+ intervals
built from cross-fitted residuals are mildly conservative at 0.907 under
representative sampling. Under the clustered design of that study they cover
0.861 against a nominal 0.90, and across the geometries we test the shortfall
runs from 0.813 to 0.883. Sample size does not repair it. Holding the primary
units fixed and quadrupling the points inside them lowers coverage further, from
0.883 to 0.856, while the map improves; adding units instead raises it. Both
axes move coverage, in opposite directions, so a larger sample is no guarantee
of a better interval. Recording every
validation point's distance to the nearest calibration point shows why: the
realised error grows with that distance while the interval width does not, and
designs that differ in marginal coverage cover almost identically at matched
distance. A spatial-block calibration holds 0.93 to 0.95 throughout, at the cost
of width.
\end{abstract}

\begin{keyword}
Spatial prediction\sep Design-based inference\sep Prediction intervals\sep
Conformal prediction\sep Variogram\sep Random forest
\end{keyword}

\end{frontmatter}

\section{Introduction}

\citet{wadoux2021} showed that map accuracy indices must be estimated from a
probability sample by design-based inference, and that spatial cross-validation
has no theoretical underpinning for this purpose. That argument concerns the
point map and its accuracy. It leaves open what should be done about the other
half of a map's quality statement, the uncertainty attached to each prediction,
and that is where this paper begins.

Prediction intervals in environmental mapping are routinely obtained from
kriging variances or from quantile heuristics, and whether they deserve to be
believed has begun to receive attention. \citet{schmidinger2023} validate
probabilistic predictions in digital soil mapping against independent test data,
scoring reliability and sharpness across a range of methods, and
\citet{rossiter2022} assess how faithfully digital soil maps reproduce the soil
geography they claim to represent. What is missing from that literature is the
step \citet{wadoux2021} took for point accuracy: treating the quantity being
validated as a parameter of the finite population and estimating it from a
probability sample by design-based inference. That is what we do here, and our
first and principal argument is that it requires no new theory, because coverage
is a population parameter of exactly the same kind as the mean squared error. Section~\ref{sec:ht} shows
that the Horvitz-Thompson estimator applies to any cell-level function, so the
apparatus \citet{wadoux2021} established for accuracy carries over to coverage
without a single new assumption, and in the case of coverage the estimator is
exactly rather than approximately unbiased. The derivation takes three lines.

The payoff is not the derivation but what it lets us see. An internal
calibration diagnostic resting on the same exchangeability assumption as the
interval cannot establish coverage when the calibration sample is not
exchangeable with the target population, and a clustered calibration sample is
exactly that case: the map must predict far from every cluster, where the
residuals it incurs are larger than any the calibration sample observed. An
external probability sample is under no such handicap, which is the whole point. Section~\ref{sec:intervals-results} measures this. Conformal intervals
that are well behaved and sharp under representative sampling cover only 0.861
against a nominal 0.90 once the calibration sample is clustered, and a
spatial-block calibration returned coverage above nominal in this experiment. The
contribution is not that coverage had never been validated, which
\citet{schmidinger2023} did, but that validating it design-based against an
independent probability sample makes the failure visible and puts the estimate
on the same footing as the accuracy indices reported beside it.

Having made the failure measurable we can ask what drives it, and the answer is
not what a practitioner would guess. Section~\ref{sec:sizes} varies the
calibration sample from 200 to 800 points along two axes and finds that
enlarging it does not repair the coverage: holding the primary units fixed and
adding points inside them makes the interval worse while making the map better,
while adding units improves it. Both axes matter and they pull in opposite
directions, so a larger calibration sample is not by itself a better one. Section~\ref{sec:distance} shows why, by recording each
validation point's distance to the nearest calibration point. The realised error
grows with that distance and the interval width does not, and designs whose
overall coverage differs by 0.07 cover almost identically at matched distance.
Clustering does not corrupt the interval so much as move the map away from the
sample that built it. For anyone planning a campaign the practical reading is
that a budget spent on more locations buys interval validity and the same budget
spent on more points per location does not.

The second argument concerns the point predictor those intervals are attached
to, and we present it as the more modest of the two. Spatial predictive mapping,
in soil science \citep{mcbratney2003} as in ecology, has increasingly favoured
machine learning over geostatistics for the prediction step,
and where geostatistics survives it usually enters by assumption: the analyst
decides residuals should be kriged, picks a variogram family by eye, fixes a
binning and an estimator. We instead place those decisions in the hyperparameter
space of the learner, selected out of fold against a null in which no correction
is applied.

We are candid about what this does and does not deliver, because an ablation
settles it. Selecting variogram parameters by cross-validation is not new; it is
long-standing practice and has been automated before \citep{li2014}, and letting
a forest represent spatial structure through distance-based features is the RFsp
idea of \citet{hengl2018}. Regression kriging and the super learner
\citep{wolpert1992,vanderlaan2007} are equally established. More to the point,
Section~\ref{sec:ablation} separates two things the phrase runs together. The
variogram-derived representation does help: removing it costs $0.172$ \Mgha{},
a fifth of the total gain, in 26 of 30 repetitions. Selecting the
variogram's parameters rather than fixing them does not: it is worth $0.008$
\Mgha{}, with a paired 95 per cent interval of $[-0.038, +0.023]$ that contains
zero, and it wins in only 8 of 30 repetitions. What it does deliver is that the null
correction is in the candidate library, so the selection never returns something
worse on its own out-of-fold criterion,
that it changes the fitted model in a little under half the repetitions, that no
anisotropic geometry survives to be selected under clustered calibration, and
that it converts an assumption into a measured statement, namely that the
conventional recipe was adequate for these data. That is worth reporting, and it is all we claim.

The two arguments are connected by construction: the same cross-fitting that
selects the correction supplies the leakage-free conformity scores from which
the intervals are built.

\section{Coverage is a population parameter, exactly as accuracy is}
\label{sec:ht}

Three distinct sampling mechanisms appear in what follows, and the argument
depends on keeping them apart. Table~\ref{tab:designs} names them. We write
\emph{calibration design} for the first and \emph{validation design} for the
third throughout, and reserve the unqualified word ``design'' for neither.

\begin{table}[!ht]
\centering
\caption{The three sampling mechanisms and what each one has to satisfy. Only
the third carries the design-based estimation argument; only the first bears on
whether the conformal guarantee applies.}
\label{tab:designs}
\small
\begin{tabular}{p{0.24\textwidth}p{0.32\textwidth}p{0.34\textwidth}}
\toprule
Component & Role & Condition it must meet \\
\midrule
Calibration sample & Fits the model and builds the intervals &
Exchangeability with the target population, for the conformal guarantee to
transfer \\
Cross-fitting folds & Produce leakage-free out-of-fold residuals and, here,
the conformity scores & Must match the interval construction they feed; see
Section~\ref{sec:intervals} \\
Validation probability sample & Estimates RMSE and population coverage &
Known, positive inclusion probabilities \\
\bottomrule
\end{tabular}
\end{table}

\subsection{Conditional and repeated-sampling targets}
\label{sec:targets}

Four quantities in this area are easy to conflate, and the contribution here
concerns only the last of them.

The first is the coverage of a realised interval map. Fix the fitted model and
the intervals it produced from one calibration sample; the proportion of the $N$
population cells whose true value falls inside its own interval is then a fixed
number, a property of that map and not of the procedure that made it.

The second averages that number over repeated calibration samples, which is a
property of the procedure.

The third is conformal marginal coverage, which is a probability over the joint
draw of the calibration sample and a new point, and is what
\citet{barber2021} bound.

The fourth is design-unbiased estimation of the first. This is the one we claim.
Conditional on the fitted map and its intervals, Eq.~(4) is exactly unbiased for
the realised coverage provided the validation sample is drawn independently of
the calibration sample and follows a probability design with known, positive
inclusion probabilities. It says nothing about whether the conformal guarantee
holds; it measures what the intervals did, whatever produced them.

The distinction matters because the two can disagree, and
Section~\ref{sec:results} is where they do. We are careful about the form of
the disagreement. CV+ guarantees $1-2\alpha$, not $1-\alpha$, so at
$\alpha=0.10$ its floor is 0.80 less a finite-sample deduction, and every
coverage we measure clears that floor: 0.861 under the clustered design of
\citet{wadoux2021}, 0.813 at the worst geometry we test, 0.811 in the
simulation. The theorem is not violated anywhere in this paper, and we make no
such claim. What fails is the thing a practitioner actually relies on. Under
representative sampling these intervals sit near the nominal 0.90 they were
asked for, comfortably above the guarantee, and it is that empirical margin,
not the theorem, that clustering consumes. A design-based estimate is what
distinguishes the two, because a check resting on exchangeability cannot see a
margin disappear when exchangeability is what has gone.



Let the domain be a finite population of $N$ cells, $U=\{1,\dots,N\}$, with true
values $z(s_i)$, and let the map be a function $\hat z: U \rightarrow
\mathbb{R}$. Once produced, the map is fixed and so is the true field. The only
randomness admitted is that of the validation sampling design, which is what
makes the assessment free of assumptions about autocorrelation or stationarity
\citep{degruijter2006,brus2011}.

Let $g: U \rightarrow \mathbb{R}$ be any cell-level function with population
mean $\bar G = N^{-1}\sum_{i \in U} g(i)$, and draw a probability sample
$S \subset U$ of size $n$ with inclusion probabilities $\pi_i$. The
Horvitz-Thompson estimator \citep{horvitz1952},
\[
\widehat{\bar G} = \frac{1}{N}\sum_{i \in S} \frac{g(i)}{\pi_i},
\eqnum{1}
\]
is design-unbiased, and the proof is one line: with $I_i = \mathbf{1}\{i \in S\}$
we have $\mathbb{E}[I_i] = \pi_i$ over repeated sampling, so
\[
\mathbb{E}\big[\widehat{\bar G}\big]
= \frac{1}{N}\sum_{i \in U} \frac{\mathbb{E}[I_i]}{\pi_i}\,g(i)
= \frac{1}{N}\sum_{i \in U} g(i) = \bar G .
\eqnum{2}
\]
Under simple random sampling without replacement $\pi_i = n/N$ and Eq.~(1)
collapses to the sample mean. Everything reported below uses a simple random
validation sample, so from here on the estimators are written in that form. For
a design with unequal inclusion probabilities, or for cluster sampling, the
$\pi$-weighted form of Eq.~(1) must be used instead, and its variance depends on
first and second order inclusion probabilities rather than on the simple
expressions given below.

The point of writing it this generally is that everything else follows from
choosing $g$. Taking $g(i) = (\hat z(s_i)-z(s_i))^2$ gives the population mean
squared error and its design-unbiased estimator,
\[
\widehat{\mathrm{MSE}} = \frac{1}{n}\sum_{i \in S}\big(\hat z(s_i)-z(s_i)\big)^2 ,
\eqnum{3}
\]
which is the estimator \citet{wadoux2021} advocate. We report
$\widehat{\mathrm{RMSE}} = \sqrt{\widehat{\mathrm{MSE}}}$ and note, since it is
usually left unsaid, that the square root is concave, so Jensen's inequality
gives $\mathbb{E}[\widehat{\mathrm{RMSE}}] \le \mathrm{RMSE}$: the root mean
squared error is not exactly design-unbiased even though the mean squared error
is. The bias is $O(n^{-1})$ and negligible here, and it is a property of the
square root rather than of the design. The mean error, obtained with
$g(i) = \hat z(s_i)-z(s_i)$, is exactly unbiased. We also report the model efficiency coefficient \citep{nash1970},
$\mathrm{MEC} = 1 - \sum_{i \in S}(\hat z_i - z_i)^2 / \sum_{i \in S}(z_i - \bar z_S)^2$.
Being a ratio of two estimated quantities it is not a design-unbiased estimator
of its population counterpart, and we report it for comparability with the
reference study rather than as an unbiased quantity.

Now take $[\hat \ell(s), \hat u(s)]$ to be a nominal $1-\alpha$ prediction
interval and set $g(i) = \mathbf{1}\{\hat \ell(s_i) \le z(s_i) \le \hat u(s_i)\}$.
Its population mean is the prediction-interval coverage probability (PICP), and
Eq.~(1) supplies the estimator at once,
\[
\widehat{\mathrm{PICP}} = \frac{1}{n}\sum_{i \in S}
\mathbf{1}\{\hat \ell(s_i) \le z(s_i) \le \hat u(s_i)\}.
\eqnum{4}
\]
Because $g$ is an indicator, $\widehat{\mathrm{PICP}}$ is a sample proportion.
Under the simple random validation sample used here it is exactly
design-unbiased, with design variance $\mathrm{PICP}(1-\mathrm{PICP})/n$ up to
the finite-population correction. Both statements are properties of this design;
under a general probability design the $\pi$-weighted estimator applies and the
variance follows the design rather than this formula. Among
valid intervals the better one is narrower, measured by the mean width
$\mathrm{MPIW}$, but width without validity is meaningless and the two must be
read together.

We stress how little this required. No new theory was needed, only the
observation that coverage is a population mean. The consequence is nonetheless
substantial: coverage estimated this way is model-free and can therefore
identify undercoverage in
an interval method without sharing any of its assumptions.

\section{A correction should enter only on out-of-fold evidence}
\label{sec:correction}

Write the base random forest \citep{breiman2001} as $\hat f$ and decompose
$z(s) = \hat f(s) + r(s)$. The residual field is not noise: a forest fitted on a
finite sample cannot represent every relation, and what it missed remains in
$r$. If $r$ is predictable then a second model $\hat g$ carries information the
map is discarding, which is the logic of regression kriging \citep{hengl2018}.

The question is how much of it to add. For $\hat y = \hat f + \lambda \hat g$
the population risk is a convex parabola in $\lambda$,
\[
\mathcal{R}(\lambda)
= \mathbb{E}[r^2] - 2\lambda\,\mathbb{E}[r\hat g] + \lambda^2\,\mathbb{E}[\hat g^2],
\eqnum{5}
\]
minimised at
\[
\lambda^\star = \frac{\mathbb{E}[r\hat g]}{\mathbb{E}[\hat g^2]},
\qquad
\mathcal{R}(\lambda^\star) = \mathbb{E}[r^2]
- \frac{\mathbb{E}[r\hat g]^2}{\mathbb{E}[\hat g^2]} \le \mathcal{R}(0).
\eqnum{6}
\]
Three things follow. At the optimum the correction cannot raise the out-of-fold
criterion the selection is run on, and lowers it in
proportion to how well $\hat g$ correlates with the residual. Adding the raw
prediction, $\lambda = 1$, is optimal only if $\hat g$ is perfectly calibrated;
a noisy $\hat g$ inflates $\mathbb{E}[\hat g^2]$ and pushes $\lambda^\star$
below one, so shrinkage is the solution of the risk minimisation and not a
safety device bolted on afterwards. And solving
$\mathcal{R}(\lambda) > \mathcal{R}(0)$ shows the correction damages the map
whenever
\[
\lambda > 2\lambda^\star .
\eqnum{7}
\]
Over-correction is therefore a real failure mode, which is what makes the next
point necessary rather than merely prudent.

Estimating $\lambda$ on the data used to fit $\hat g$ fails, and fails in a
predictable direction. A forest interpolates its training data to some degree,
so the in-sample residual $y_i - \hat f(x_i)$ is shrunk toward zero and is
correlated with the noise in $y_i$. A model fitted to those residuals partly
relearns that noise and appears to explain more of $r$ than it can out of
sample, so the selected $\lambda$ is biased upward, which is exactly the regime
of Eq.~(7). The remedy is $K$-fold cross-fitting: fit the base model on
$U \setminus F_k$ and predict $F_k$ to obtain $\tilde f_i$, and likewise obtain
$\tilde g_i$ from the training folds' residuals. Neither depends on $y_i$, so
\[
\widehat{\mathcal{R}}(c,\lambda)
= \sqrt{\frac{1}{n}\sum_{i=1}^{n}\big(y_i - \tilde f_i - \lambda\,\tilde g_i^{(c)}\big)^2}
\eqnum{8}
\]
estimates out-of-sample risk rather than in-sample fit, the same device that
underlies double machine learning \citep{chernozhukov2018}.

The correction is a library of candidates: a residual forest, a residual Cubist
\citep{quinlan1992}, a single-hidden-layer network, ordinary kriging of the
residuals, a radial basis or kernel smoother, their ensemble, and linear or
quadratic recalibration. We take $(c^\star,\lambda^\star)$ to minimise Eq.~(8).
The structural detail that matters is that $\lambda = 0$ belongs to the grid and
reproduces the base forest, so
\[
\widehat{\mathcal{R}}(c^\star,\lambda^\star) \le \widehat{\mathcal{R}}(\cdot\,,0)
\eqnum{9}
\]
by construction, on the criterion the selection uses: the method cannot select
something that scores worse out of fold than the plain forest
on the out-of-fold criterion, and a correction is deployed only after showing an
improvement on data it did not see. This is why the estimator degrades to a
random forest, rather than to something worse, when there is nothing to recover.

Two qualifications belong in the open. The minimum of $\widehat{\mathcal{R}}$
over many candidates is optimistically biased, the familiar winner's curse, so
it is not an unbiased estimate of deployed risk and we never report it as one.
That is precisely why selection happens inside the calibration sample while
accuracy is computed on an independent probability sample: the winner's curse
contaminates the selection score, not the design-based estimate.

Finally, a correction is credible only where data support it. With $d_1(s)$ the
distance to the nearest calibration point we gate it,
\[
w(s) = \exp\!\big(-d_1(s)/\ell\big),
\qquad
\hat y(s) = \hat f(s) + \lambda^\star\, w(s)\, \hat g^{(c^\star)}(s),
\eqnum{10}
\]
with $\ell$ tied to the fitted variogram range, since beyond that range
residuals are effectively uncorrelated and an interpolated correction carries no
information about them. The multiplier is itself selected out of fold, including
the value that disables the gate.

\section{The variogram is a set of hyperparameters, not a set of conventions}
\label{sec:geohyper}

The quantities above are selected, but the variogram feeding them is normally a
fixed recipe: a family chosen by habit, a cutoff at half the maximum distance,
a fixed lag count, one estimator, and a range consumed as a constant. For an
argument that geostatistical choices should rest on evidence, that is the last
unexamined assumption, and we remove it.

The family is selected rather than assumed. With the empirical variogram
computed by the robust estimator of \citet{cressie1980}, the exponential,
spherical and Gaussian models \citep{matheron1963,journel1978}, together with
the Mat\'ern model \citep{matern1960,stein1999}, are
each fitted to the same empirical variogram and the lowest weighted residual sum
of squares is retained. The comparison is legitimate in this direction, and only
in this direction, because binning and pair counts are identical across
candidates, so the criteria are computed on the same quantity. Parsimony is
enforced among nested families: Mat\'ern nests the exponential at smoothness
$\kappa = 1/2$ and approaches the Gaussian as $\kappa \rightarrow \infty$, and
the smoothness is itself fitted, so raw sum of squares hands Mat\'ern a win by
one extra degree of freedom even when the truth is exactly exponential, which we
observed. We therefore keep the simplest family within a tolerance of the best.

The fitted range enters as a working scale rather than a constant, with radial
basis bandwidth $b = m_b\,\rho$ and the multiplier selected out of fold; the
conventional value corresponds to $m_b = 1/2$. The lag binning and the choice
between the robust and the classical \citep{matheron1963} estimator cannot be
compared by variogram sum of squares at all, because changing them changes the
empirical variogram and hence the quantity being fitted. They are selected by
out-of-fold predictive merit through Eq.~(8).

Anisotropy is treated the same way. We estimate a major-axis angle $\theta$ and
a ratio $\varrho \in (0,1]$ and map coordinates into the space where Euclidean
distance is anisotropic distance,
\[
T_{\theta,\varrho}(s) =
\begin{pmatrix}
\sin\theta & \cos\theta \\[2pt]
-\cos\theta/\varrho & \sin\theta/\varrho
\end{pmatrix} s ,
\eqnum{11}
\]
so that $\varrho = 1$ is a pure rotation recovering the isotropic method
exactly. The transform is applied to the radial basis, the kernel smoother and
the kriging correction, since kriging under geometric anisotropy is ordinary
kriging in the transformed space \citep{journel1978}.

Estimating $(\theta,\varrho)$ deserves a warning, because the obvious routes
fail and we tried them. Taking the ratio of the smallest to the largest
directional range is biased by order statistics, since the maximum of noisy
estimates overshoots and the minimum undershoots; it reports anisotropy for
perfectly isotropic fields. Scoring geometries by variogram fit in the
transformed space fails differently, because the transform changes the distance
distribution and hence the binning, so the criteria are not comparable. We
therefore score each candidate geometry by out-of-fold prediction error of a
kernel smoother on the same held-out points, with bandwidth set as a fraction of
the median pairwise distance in the transformed space so that the score responds
to shape and not to rescaling. On simulated fields drawn by simple random
sampling this returns exact isotropy for isotropic fields and the correct angle
otherwise, and its errors there are false negatives; under clustered sampling
specificity is high but not perfect (Section~\ref{sec:sims}). Even then the geometry is not imposed: the anisotropic variant enters
the search as one more candidate and must beat the isotropic one on Eq.~(8).

\section{Conformal intervals, and why clustering breaks them}
\label{sec:intervals}

Cross-fitting has already produced leakage-free residuals for the deployed
configuration, and these are exactly the conformity scores conformal prediction
requires \citep{papadopoulos2002,vovk2005,lei2018}.

The argument is clearest in the split conformal case, where one calibration set
is held out. With $q$ the $\lceil (n+1)(1-\alpha)\rceil$-th smallest $|e_i|$,
each cell receives $[\hat y(s)-q,\ \hat y(s)+q]$, and the guarantee rests on
exchangeability rather than on distributional form. If
$|e_1|,\dots,|e_n|,|e_{\mathrm{new}}|$ are exchangeable then $|e_{\mathrm{new}}|$
is equally likely to occupy any rank among the $n+1$ values, so
\[
\Pr\big(|e_{\mathrm{new}}| \le q\big) \ge 1-\alpha
\eqnum{12}
\]
with no assumption of Gaussianity, homoscedasticity or stationarity
\citep{vovk2005}.

The construction is worth stating step by step, because the words
``cross-fitted residuals'' cover several inequivalent procedures. Write
$\mathcal{S}$ for the calibration sample, partitioned into $K$ folds
$\mathcal{S}_1,\dots,\mathcal{S}_K$.

\begin{enumerate}\itemsep2pt
\item For each fold $k$, run the \emph{entire} selection, the profile search,
the correction mode, the shrinkage $\lambda$, the gate and the variogram
recipe, using only $\mathcal{S}\setminus\mathcal{S}_k$. Call the resulting
fitted predictor $\hat f^{-k}$.
\item For every calibration point $i\in\mathcal{S}_k$ compute
$r_i = |y_i - \hat f^{-k}(x_i)|$.
\item For every target cell $s$ compute $\hat f^{-k}(s)$, using the same
fold-$k$ model.
\item Form the $n$ candidate lower bounds $\hat f^{-k(i)}(s) - r_i$ and the $n$
candidate upper bounds $\hat f^{-k(i)}(s) + r_i$, where $k(i)$ is the fold
containing $i$.
\item Take the $\lfloor \alpha(n+1)\rfloor$-th smallest lower bound and the
$\lceil (1-\alpha)(n+1)\rceil$-th smallest upper bound
\citep{barber2021}.
\end{enumerate}

Two features distinguish this from the alternatives. The interval is not
centred on a single prediction: it aggregates bounds from $K$ different models,
so it is asymmetric in general. And the selection is repeated inside every fold
rather than performed once on all of $\mathcal{S}$. The second point is not
cosmetic. Choosing the configuration on the same out-of-fold residuals that
later serve as conformity scores would let a configuration be selected
\emph{because} its residuals were small, and those residuals would then no
longer be exchangeable with a fresh one. That is the selection-induced optimism
we acknowledge for risk estimation in Section~\ref{sec:correction}, and it applies
to calibration too. Nesting the selection costs $K$ times the search, and it is
what makes the bound in Eq.~(12) applicable to the interval we actually report.

The point map is a separate object. It comes from the configuration selected on
all of $\mathcal{S}$ and refitted there, as Fig.~\ref{fig:schematic} shows. The
interval is not a band around that map.

We do not deploy that construction. Holding out a separate calibration set is
costly at the sample sizes environmental mapping campaigns work with, so the
intervals
reported here are CV+ \citep{barber2021}, formed over the same $K$ cross-fitting
folds and reusing every observation for both fitting and calibration. The price
is a weaker guarantee, still finite-sample and still resting only on
exchangeability: CV+ attains at least $1-2\alpha$, not the $1-\alpha$ of
Eq.~(12). That bound is conservative in practice, and
Section~\ref{sec:sims} measures how much.

When calibration and target cells are drawn by simple random sampling from the
same population, exchangeability holds and the coverage statement is one about
the population, which is design-based coverage in the sense of Eq.~(4).

Simple random sampling of the calibration points makes them exchangeable with a
uniformly drawn target cell in the sense the guarantee needs: sampling without
replacement from a finite population gives symmetry under permutation, not
independence, and permutation symmetry is what the rank argument uses.

Two exactness claims appear in this paper and they rest on different things, so
we separate them. The design-based coverage estimator of Eq.~(4) is exactly
unbiased for any probability design with known positive inclusion
probabilities. That is a property of the estimator and the validation design
alone; it holds whatever produced the intervals, and nothing about
exchangeability enters it. The conformal bound of Eq.~(12) is the claim that
needs exchangeability, and it is exact when the calibration points and the
target form one simple random sample of $n+1$ cells drawn without replacement.
Our validation sample is drawn independently of the calibration sample from the
same frame and is not screened against it, so the two are not quite that object:
a validation cell coincides with a calibration cell with probability $n/N$,
here $3.8\times10^{-4}$, which puts an expected 0.19 coincident cells in a
500-cell validation draw and leaves about one repetition in six with a single
one. The departure from exchangeability is of that order and we have not
corrected for it. Screening the validation draw against the calibration sample
would remove it, at the cost of a validation design that is no longer a simple
random sample of the population the map is meant to cover, which is the estimand
Section~\ref{sec:ht} defines. Our
systematic design does not have that property. It draws a randomly located
regular grid, so the design is probabilistic and the validation estimator is
unbiased, but the calibration points are not exchangeable with an arbitrary
target and Eq.~(12) does not transfer to it. The coverage we report for the
systematic design is therefore an empirical finding, not a guarantee attained.


Clustered sampling breaks it, and does so directionally. The map must predict
everywhere, including far from every cluster, and the residuals it incurs there
are systematically larger than any observed within clusters, because those
locations are extrapolations. The calibration scores are concentrated on
well-supported locations, so $q$ is calibrated on residuals that are too small,
the intervals are too narrow, and coverage falls below nominal. We test two
repairs. Covariate-shift weighting replaces the empirical quantile by a weighted
one with density-ratio weights \citep{tibshirani2019}; its limitation is
structural rather than technical, since reweighting can only redistribute mass
over residuals that were actually observed. Spatial-block conformal calibration
instead changes what the scores are, refitting on spatial blocks and scoring each
held-out block \citep{roberts2017}, so that every residual is generated in an
extrapolation regime resembling the one the map faces.

\section{Experimental design}
\label{sec:case}

We use the public materials of \citet{wadoux2021}: above-ground live woody
biomass surface of \citet{baccini2012} over a 928 by 1642 km Amazonian rectangle at 1 km
resolution, 1{,}310{,}143 cells, with 28 covariates known over the whole grid.
Following \citet{wadoux2021} we treat that surface as the finite population for
the numerical experiment. We frame the paper as spatial predictive mapping
rather than digital soil mapping specifically: the argument is about how any
map built from a point sample of a spatial variable should have its intervals
assessed, and the benchmark is above-ground biomass rather than a soil
property. Nothing in the derivation is particular to soil, but neither have we
demonstrated it on a pedological survey, and Section~\ref{sec:limits} says so. It is itself a remote-sensing estimate carrying its own
error rather than field-measured biomass, so the results below concern how well a
mapping procedure reconstructs this benchmark pseudo-population.
Calibration samples of $n = 500$ are drawn under the same three designs they
used: simple random, systematic on a randomly located grid, and two-stage
clustered random, the last drawing 20 primary units and 25 points within each so
that the sample covers a small fraction of the domain and most of the map is an
extrapolation. Each map is evaluated on an independent simple random validation
sample of 500 cells, drawn independently of the calibration sample, with 30
repetitions
per design and a fixed seed. Calibration and validation sizes are therefore both
$n = 500$ for these three headline campaigns. Two further sets are introduced
where they are used: Section~\ref{sec:sizes} sweeps the calibration size from
200 to 800, on the clustered design along two axes, the number of primary units
and the points drawn within each, with the validation draw held at 500
throughout and 30 repetitions per cell; and Section~\ref{sec:distance} re-runs
three of those designs for 10 repetitions each, recording every validation
point's distance to the nearest calibration point together with the interval it
received. Comparators are the untuned default
forest exactly as in the reference study, a default Cubist, and a transparency
model we label RF-tuned, the selected base forest with no correction applied,
which separates the effect of tuning the forest from the effect of the
correction. Forests use \texttt{ranger} \citep{wright2017} and all variogram
work uses \texttt{gstat} \citep{pebesma2004}; every forest receives an explicit
seed. The estimator is released as the \texttt{GeoVersa} package for R.
Fig.~\ref{fig:schematic} summarises the procedure. We refer to the predictor
as TARF, a targeted adaptive residual forest, the estimator the \texttt{GeoVersa}
package names DB-TARF, and deliberately keep
``design-based'' out of that name: the learner is fitted by ordinary means,
with no Horvitz-Thompson estimator and no design-weighted objective. What is
design-based here is the assessment protocol, not the model.

\begin{figure}[!ht]
\centering
\includegraphics[width=0.8\textwidth]{figures/fig1_schematic.png}
\caption{The procedure. Cross-fitting on the calibration sample produces
leakage-free out-of-fold predictions and residual candidates; the profile,
correction, shrinkage and gate are selected by out-of-fold RMSE and accepted
only if they beat the RF-only baseline; the winners are refitted on the full
sample and the map is scored on an independent probability sample.}
\label{fig:schematic}
\end{figure}

\section{Results}
\label{sec:results}

\subsection{The correction improves the map, and the tuning does not}

The estimator lowers the root mean squared error under all three designs
(Table~\ref{tab:accuracy}), winning 28, 29 and 29 of 30 paired repetitions, the
systematic design contributing one exact tie and no losses, with two-sided
Wilcoxon $p$ values of $1.3\times10^{-8}$, $3.7\times10^{-9}$ and
$3.7\times10^{-9}$, and design-based bias moves toward zero in every design. The
RF-tuned row is the one to read carefully: tuning the forest buys between 0.008
and 0.160 \Mgha{}, one to fifteen per cent of what the full method achieves, so
the improvement is the residual correction and not a better forest. One estimand in Table~\ref{tab:estimands} is a large independent simple random
sample of 20{,}000 cells, which is a precise estimate of population accuracy but
not the population quantity itself. That quantity is available here, because the
benchmark surface is fully known: computed over all 1{,}310{,}143 cells, the root mean squared error is 31.655 \Mgha{} for the default
forest and 31.079 for the corrected map, with bias falling from $+0.378$ to
$+0.127$ (Fig.~\ref{fig:maps}).

\begin{table}[!ht]
\centering
\caption{Design-based accuracy by calibration design, 30 repetitions each,
$n=500$ for both calibration and validation. ME is the mean error and RMSE the
root mean squared error, both in \Mgha{}; SD is the standard deviation of the
RMSE across repetitions; MEC is the modelling efficiency coefficient of
\citet{nash1970}, one minus the ratio of the residual to the total sum of
squares. TARF is the method proposed here.}
\label{tab:accuracy}
\small
\begin{tabular}{llrrrr}
\toprule
Design & Model & ME & RMSE & SD(RMSE) & MEC \\
\midrule
Simple random & RF (default) & +1.061 & 32.493 & 1.219 & 0.876 \\
Simple random & RF-tuned & +1.153 & 32.466 & 1.242 & 0.875 \\
Simple random & Cubist & +1.853 & 33.264 & 1.424 & 0.870 \\
Simple random & \textbf{TARF} & \textbf{+0.802} & \textbf{31.644} & \textbf{1.307} & \textbf{0.883} \\
\midrule
Systematic & RF (default) & +0.664 & 32.280 & 1.146 & 0.877 \\
Systematic & RF-tuned & +0.731 & 32.272 & 1.146 & 0.878 \\
Systematic & Cubist & +1.921 & 33.171 & 1.152 & 0.870 \\
Systematic & \textbf{TARF} & \textbf{+0.499} & \textbf{31.400} & \textbf{1.004} & \textbf{0.884} \\
\midrule
Clustered random & RF (default) & -1.356 & 35.778 & 2.562 & 0.851 \\
Clustered random & RF-tuned & -1.220 & 35.617 & 2.555 & 0.853 \\
Clustered random & Cubist & +1.771 & 37.155 & 4.466 & 0.837 \\
Clustered random & \textbf{TARF} & \textbf{-0.876} & \textbf{34.678} & \textbf{2.506} & \textbf{0.860} \\
\bottomrule
\end{tabular}

\end{table}

The paired differences carry the comparison, since every repetition scores both
maps on the same calibration draw and the same evaluation sample.
Table~\ref{tab:paired} reports them with their dispersion, a Monte Carlo
standard error for the mean and a confidence interval, so that the reader can
separate the size of the effect from the precision with which 30 repetitions
pin it down. The interval excludes zero under all three designs, and the Monte
Carlo error is roughly a tenth of the effect, so the ranking is not an artefact
of the number of repetitions.

\begin{table}[!ht]
\centering
\caption{Paired difference in RMSE between the corrected map and the default
forest, by calibration design, over 30 repetitions. Negative values favour the
corrected map. MCSE is the Monte Carlo standard error of the mean difference,
and the confidence interval is the 95 per cent interval for that mean. RMSE in
\Mgha{}.}
\label{tab:paired}
\footnotesize
\setlength{\tabcolsep}{2pt}
\begin{tabular}{lrrrrrr}
\toprule
Design & Mean $\Delta$ & SD & MCSE & 95\% CI & Median [IQR] & W / L / T \\
\midrule
Simple random & $-0.849$ & 0.48 & 0.088 & $[-1.03, -0.67]$ & $-0.88$ $[-1.10, -0.53]$ & 28 / 2 / 0 \\
Systematic & $-0.880$ & 0.45 & 0.082 & $[-1.05, -0.71]$ & $-0.95$ $[-1.16, -0.56]$ & 29 / 0 / 1 \\
Clustered random & $-1.099$ & 0.71 & 0.129 & $[-1.36, -0.84]$ & $-1.00$ $[-1.38, -0.64]$ & 29 / 1 / 0 \\
\bottomrule
\end{tabular}

\end{table}

\begin{figure}[!ht]
\centering
\includegraphics[width=\textwidth]{figures/fig7_maps.png}
\caption{Observed biomass, the two predicted surfaces and their error fields for
one calibration sample of $n=500$. Panels (a) to (c) share one colour scale and
panels (d) and (e) share one symmetric scale, so that every comparison is visual
as well as numerical. Red circles in (a) mark the calibration points.}
\label{fig:maps}
\end{figure}

\subsection{The data displaces the conventional variogram recipe}

Table~\ref{tab:axes} reports how often the out-of-fold search departs from the
conventional recipe. It is kept in 16, 18 and 19 of the 30 repetitions, so it is
displaced in a little under half: often enough that the defaults are not
automatic, but not in a majority.

\begin{table}[!ht]
\centering
\caption{Repetitions, out of 30 per design, in which the out-of-fold search departs from the conventional recipe on each of the three axes it is recorded against: the binning and estimator (a cutoff at half the maximum distance, 15 lags, the robust estimator), the range scale (a bandwidth of half the practical range), and isotropy. The selected variogram family is not recorded per repetition and so is not counted here.}
\label{tab:axes}
\small
\setlength{\tabcolsep}{2pt}
\begin{tabular}{lrrrr}
\toprule
Design & Binning or estimator & Range scale & Anisotropy & Recipe kept \\
\midrule
Simple random & 7 & 4 & 3 & 16 \\
Systematic & 3 & 7 & 2 & 18 \\
Clustered random & 4 & 7 & 0 & 19 \\
\bottomrule
\end{tabular}

\end{table}

Two features of the clustered design deserve attention. The anisotropic geometry
is accepted in none of its 30 repetitions, and the candidate scores confirm that
it was never in contention: the anisotropic profile and its isotropic twin
record identical out-of-fold errors in all 30, so the geometry estimator's 3 per
cent parsimony threshold settled the matter before the search saw it. The range
scale is the axis the search displaces most often, in 7. The binning and estimator move in 4, fewer than
under simple random sampling, where they move in 7.

We resist reading a mechanism into the anisotropy result at this point. The
natural explanation, that a clustered sample carries too little information
about direction to support a geometry, is one we test directly in
Section~\ref{sec:sims} and fail to confirm. What the counts do show is that the
search does not behave as a switch that turns geostatistics on: it leaves the
conventional recipe untouched in 19 of the 30 clustered repetitions, and moves
individual axes only where the sample gives it reason to.

Table~\ref{tab:modes} reports which corrector is selected across all 90
repetitions, and it is also the limit of what we claim. Ordinary kriging wins
outright in three repetitions and a purely geostatistical interpolator in four;
in 72 the corrector is a residual forest which, carrying a shrinkage weight, is
formally a single boosting step. It is fair to ask whether this is
geostatistics at all.

\begin{table}[!ht]
\centering
\caption{Selected corrector across all 90 repetitions.}
\label{tab:modes}
\small
\begin{tabular}{lr}
\toprule
Selected corrector & Count \\
\midrule
Residual random forest & 72 \\
Ensemble of residual learners & 10 \\
Ordinary kriging of residuals & 3 \\
Cubist delta & 1 \\
Linear recalibration & 1 \\
Residual Cubist & 1 \\
Single-hidden-layer network & 1 \\
Kernel (radial basis) smoother & 1 \\
\bottomrule
\end{tabular}

\end{table}

It is, but in the representation rather than in the estimator. The residual
forest is not fitted on coordinates alone: its features include radial basis
functions whose centres and bandwidth come from the fitted variogram, and from
the anisotropic geometry of Eq.~(11) when that is accepted, and its output is
multiplied, where the search retains it, by the gate of Eq.~(10) whose length
scale is the fitted range. The search disabled that gate in 81 of the 90
repetitions, keeping it only under clustered calibration, so the features carry
the argument here and the gate does not. Remove the variogram and neither
exists. What
Table~\ref{tab:modes} shows is that once spatial structure is encoded this way a
forest exploits it better than a kriging predictor does on these data, which is
a statement about two learners and not evidence that the geostatistical
information was unnecessary.

\subsection{The advantage is not specific to the scoring protocol}

We drew 10 further calibration samples under simple random sampling and scored
each under three additional protocols (Table~\ref{tab:estimands}): a Population
estimand, an independent simple random sample of 20{,}000 cells excluding the
calibration cells; random 10-fold cross-validation; and spatial block
cross-validation at 350 km, a distance chosen because it exceeds the practical
range fitted to these residuals, so that held-out blocks fall beyond the
correlation the model can exploit.

Two qualifications bound what this shows, and following \citet{wadoux2021} we
state them plainly. Only the first two rows are map-accuracy estimands; the
cross-validation protocols are not valid estimators of map accuracy and are
reported for model comparison only. They are, moreover, aligned with the
internal selection criteria of the method, since corrections are accepted on
random-fold out-of-fold risk and the gate is chosen by spatial
cross-validation, which makes them the least independent of the four checks
rather than the most adversarial. The Population row carries the evidential
weight.

\begin{table}[!ht]
\centering
\caption{Comparison against the untuned default forest under four estimands,
simple random calibration design. RMSE in \Mgha{}.}
\label{tab:estimands}
\small
\begin{tabular}{lrrrrr}
\toprule
Estimand & Reps & RF & TARF & $\Delta$RMSE & Wins \\
\midrule
\multicolumn{6}{l}{\emph{Design-based map accuracy, independent probability samples}} \\
Large independent SRS (20{,}000) & 10 & 32.17 & 31.51 & -0.669 & 10/10 \\
Design-based (SRS 500) & 30 & 32.49 & 31.64 & -0.849 & 28/30 \\
\midrule
\multicolumn{6}{l}{\emph{Model comparison only, not map accuracy}} \\
Random 10-fold CV & 10 & 32.24 & 31.11 & -1.125 & 10/10 \\
Spatial block CV (350 km) & 10 & 34.72 & 33.45 & -1.275 & 10/10 \\
\bottomrule
\end{tabular}

\end{table}

\begin{figure}[!ht]
\centering
\includegraphics[width=0.88\textwidth]{figures/fig4_protocols.png}
\caption{Comparison under four validation estimands, simple random calibration
design. Population and Design-based estimate map accuracy; the two
cross-validation estimands are reported for model comparison only and are not
valid map-accuracy estimates. The cross-validation bars train on subsets of the
calibration sample, so their absolute heights are not comparable with the first
two.}
\label{fig:protocols}
\end{figure}

The method wins all 10 draws under each additional protocol
(Fig.~\ref{fig:protocols}). Three cautions. The
one-sided paired Wilcoxon $p$ at 10 repetitions cannot fall below
$2^{-10} \approx 0.00098$, the value attained, so it certifies that all ten
draws favour the method and nothing stronger. These are 10 draws scored three
ways rather than 30 independent repetitions, and results within a draw are
dependent, the strongest of the three pairwise correlations between protocols
reaching $r = 0.59$. The cross-validation rows train on strict
subsets, so their absolute values are not commensurable with the design-based
rows. We do not read the ordering of the four differences as a gradient:
adjacent protocols are not statistically distinguishable ($p = 0.36$), and the
ordering is not one of spatial stress, since random 10-fold cross-validation is
the least spatially demanding of the four. Buffered leave-one-out
cross-validation was not run, requiring of order $10^3$ sequential fits per
repetition; it would be the most adversarial estimand available and is left to
future work.

\subsection{Design-based assessment identifies the undercoverage}
\label{sec:intervals-results}

Table~\ref{tab:coverage} and Fig.~\ref{fig:uncertainty} contain the finding we
consider most consequential.
Under representative sampling the plain conformal intervals are calibrated,
0.907 and 0.913 against a nominal 0.90, mildly conservative in the direction
CV+ predicts, and about 16 per cent narrower than the quantile regression forest
baseline \citep{meinshausen2006}. That baseline over-covers at 0.937, so the two
are not compared at the same empirical coverage and the width difference is
descriptive: a method that covers less will generally be narrower. A comparison
at matched coverage would need both recalibrated to a common level, which we did
not do. Under the clustered design of \citet{wadoux2021}, twenty primary units
of twenty-five points each, they under-cover at 0.861, with a 95 per cent
confidence interval for the mean of $[0.846, 0.875]$ that excludes the nominal
level, and 25 of the 30 repetitions falling below it individually. That figure
belongs to that geometry and not to clustered sampling in general:
Section~\ref{sec:sizes} varies the number of primary units and the points within
them and finds coverage ranging from 0.813 to 0.883 across the geometries
tested, all of them short of nominal.

Two features of that failure are worth separating, and they are visible in
Fig.~\ref{fig:uncertainty} rather than in any mean. The first is the shift: the
whole distribution moves down, so this is not a handful of bad repetitions
dragging an otherwise sound average. The second is the spread. The
between-repetition standard deviation roughly doubles under clustering, from
$0.019$ to $0.040$, so a clustered calibration sample makes coverage both lower
and less predictable. A practitioner drawing one such sample cannot tell from
the sample itself which end of that distribution they landed on.

Covariate-shift weighting does not repair this: at 0.862, with 26 of 30
repetitions below nominal, it is indistinguishable from the unweighted interval, which is
what the structural argument of Section~\ref{sec:intervals} predicts, since
reweighting redistributes mass over residuals the calibration sample never
observed. The spatial-block calibration recovered conservative empirical
coverage in this experiment, reaching 0.946, above the nominal 0.90, at
the cost of intervals about 48 per cent wider.

The quantile regression forest deserves more attention here than a baseline
usually gets, because it does not fail. Its coverage is 0.937, 0.936 and 0.928
across the three designs, so clustering costs it 0.009 where it costs CV+
0.046, and under clustered calibration it reaches 0.928 with intervals of 137
\Mgha{} against the spatial-block repair's 0.946 at 144. We report this because
it is what the experiment found and because it qualifies the framing: a
practitioner whose only concern is that clustered intervals should not
under-cover could reach for the quantile forest and be served. What that method
does not offer is the reason to expect it. It carries no finite-sample bound, it
over-covers by 0.037 under representative sampling where CV+ is calibrated to
0.007, and it is wider than CV+ everywhere, so its robustness under clustering
is bought in advance and paid for on every map. Its widths are also per point,
derived from the spread of the forest's own leaves, which is a plausible reason
it degrades less where the model extrapolates; our per-point record covers only
the conformal interval, so we cannot test that here and offer it as a
conjecture rather than a finding.

\begin{table}[!ht]
\centering
\caption{Design-based interval coverage at a nominal 0.90 level and mean width in
\Mgha{}. PICP and MPIW are means over 30 repetitions, with leading zeros dropped on
the bounded columns; SD is the standard
deviation of PICP across them and the interval is a 95 per cent confidence
interval for that mean, from the Monte Carlo standard error. The last column
counts repetitions whose own coverage fell below nominal. OOF is out of fold
and RF is random forest. The rows are not all the same construction, and the comparison is
empirical rather than like-for-like. CV+ is the deployed interval and the only
row carrying a finite-sample bound. The pooled out-of-fold quantile centres a
single half-width on the full-sample refit and is reported because it is what a
natural reading of cross-fitted residuals produces, not because it is valid. The
covariate-shift and spatial-block rows modify that pooled calibration, by
reweighting and by rescoring on held-out blocks respectively.}
\label{tab:coverage}
\footnotesize
\setlength{\tabcolsep}{2pt}
\begin{tabular}{llrr}
\toprule
Design & Interval & PICP & MPIW \\
\midrule
\multirow{2}{*}{Simple random} & Conformal & 0.895 & 100 \\
 & Quantile regression forest & 0.937 & 123 \\
\midrule
Systematic & Conformal & 0.902 & 103 \\
\midrule
\multirow{3}{*}{Clustered random} & Conformal, plain & 0.849 & 94 \\
 & Covariate-shift weighted & 0.863 & 98 \\
 & \textbf{Spatial-block conformal} & \textbf{0.946} & \textbf{144} \\
\bottomrule
\end{tabular}

\end{table}

What matters is not any single number but that the failure is visible at all. A
model-based coverage check would have relied on the same exchangeability that
clustering breaks, and would have reported these intervals as adequate.
Estimating coverage as a population parameter on a probability sample is what
identifies it, and it does so without sharing any assumption with the method
under test.

\begin{figure}[!ht]
\centering
\includegraphics[width=\textwidth]{figures/fig6_uncertainty.png}
\caption{Design-based interval assessment. Left, coverage with the nominal 0.90
level marked; right, mean interval width. The plain conformal calibration is
valid and sharpest under representative sampling but under-covers under
clustering. Two intervals hold up there, the spatial-block calibration and the
quantile regression forest, both at the cost of width.}
\label{fig:uncertainty}
\end{figure}

\subsection{The failure is indexed by distance from the calibration sample}
\label{sec:distance}

Sections~\ref{sec:intervals} and \ref{sec:intervals-results} assert that the
mechanism is spatial extrapolation but measure coverage only over the whole map.
Fig.~\ref{fig:distance} measures it against the variable the mechanism names. We
re-ran three designs recording, for every validation point, its distance to the
nearest calibration point together with the interval it received: 10 repetitions
each of simple random calibration at $n = 500$, the clustered design of
\citet{wadoux2021} at the same size, and the eight-unit geometry, which carries
200 calibration points rather than 500 and so differs from the others on both
axes at once; 5000 validation points per design.

Coverage declines with that distance under all three. Fitting one logistic slope
per repetition and testing the ten slopes, a miss becomes more likely by
$+0.189$ log-odds per doubling of distance under simple random calibration
($[+0.139, +0.240]$), $+0.198$ under twenty units ($[+0.104, +0.292]$) and
$+0.363$ under eight ($[+0.244, +0.483]$). The within-repetition standard errors
play no part in those intervals; only the spread between repetitions does, so
nothing here treats points that share a fitted model as independent.

Two things are worth separating. The first is that the three designs behave
almost identically at matched distance, which Fig.~\ref{fig:distance}c shows and
which the marginal coverages conceal. Comparing each clustered design with
simple random inside the four distance bands they share, seven of the eight
contrasts contain zero, and the one that does not is a shortfall of $0.035$
between 25 and 50 km for the twenty-unit design, with a Holm-adjusted
bootstrap $p$ near $0.02$. What differs between designs is not how they behave at a given
distance but how much of the map sits at each distance. The median distance from
a validation point to the nearest calibration point is 25 km under
simple random sampling, 75 under twenty units and 164 under eight, and the
fraction of the domain lying more than 100 km from any calibration
point is 0.001, 0.387 and 0.687. Clustering does not corrupt the interval so
much as move the population away from the sample that built it.

The second is what the interval does about that, which is nothing.
Fig.~\ref{fig:distance}b puts the realised error and the issued half-width on
one axis. Under the eight-unit geometry the conditional RMSE rises from 26.8 to
48.4 \Mgha{} between the nearest and farthest bands while the half-width moves
from 51.5 to 53.8, a rise of 81 per cent against 5 per cent. The CV+ construction
emits a separate bound for every point and those bounds do vary, with a standard
deviation of 2.2 \Mgha{} within a repetition of that design, so the width is
free to respond to location. It does not respond to distance, because the conformity scores that
set it were collected where the calibration sample is.

The residual band contrast has the same origin. At matched distance the
twenty-unit design issues a narrower interval than simple random, 48.6 against
51.2 \Mgha{}, and its out-of-fold error is correspondingly smaller, 29.4 against
31.0, both read from the same ten repetitions as the rest of this section. A sample whose points sit near one another is easy to interpolate between,
so it looks more predictable than the domain is, and the conformal quantile
inherits the optimism. Clustering therefore acts twice and in the same
direction: it moves the target population away from the calibration sample, and
it makes that sample understate the error.

\begin{figure}[!ht]
\centering
\includegraphics[width=\textwidth]{figures/fig8_distance.png}
\caption{Where the interval fails. (a) The fraction of the mapped domain lying
farther than $\delta$ from any calibration point, estimated design-unbiasedly
because validation is a simple random sample of the same frame in every
repetition. (b) Conditional RMSE (solid) and the mean half-width the interval
issued (dashed), in distance bands. (c) Coverage within each band, with 95 per
cent intervals from a bootstrap that resamples whole repetitions, since the 500
points of a repetition share one fitted model and one conformal calibration; the
faint lines are the coverage of all points beyond $\delta$, which cannot
coincide across designs because each averages over its own distance
distribution. Distance indexes the failure and is not shown to cause it: within
one design, distance from the calibration sample and geographic position are the
same variable, and the causal argument is the stationary-field simulation of
Section~\ref{sec:sims}. Distances are geometric and say nothing directly about
support in covariate space.}
\label{fig:distance}
\end{figure}

\subsection{More data does not repair the coverage, and can worsen it}
\label{sec:sizes}

The obvious reading of Section~\ref{sec:intervals-results} is that a clustered
sample of 500 points is too small, and that a larger one would supply the
residuals the conformity scores are missing. Table~\ref{tab:sizes} tests it by
varying the calibration size from 200 to 800 along two axes, 30 repetitions per
cell, with the validation draw held at 500 throughout so that every coverage
estimate is measured at the same precision.

Hold the geometry fixed at twenty primary units and add points inside them, and
coverage falls. It is 0.883 at $n=200$ and 0.856 at $n=800$, a change of
$-0.027$ with a 95 per cent confidence interval of $[-0.044, -0.010]$, while
RMSE improves from 36.85 to 34.22 \Mgha{}. Quadrupling the calibration sample
therefore bought a better map and a worse interval. We do not read the three
means as a monotone sequence: the step from 500 to 800 is $-0.005$ and is not
resolved by 30 repetitions. The end-to-end contrast is what the data support.

Hold the points per unit fixed and add units instead, and coverage rises, from
0.813 to 0.880, a change of $+0.067$ with interval $[+0.044, +0.089]$. The two
axes can be compared at equal cost. At $n=200$, spending the budget on twenty
units rather than eight is worth $+0.069$ with interval $[+0.046, +0.093]$; at
$n=800$, thirty-two units rather than twenty is worth $+0.024$ with interval
$[+0.009, +0.039]$. Both survive a Holm correction over the eight contrasts
these campaigns admit, adjusted $p<10^{-4}$ and $p=0.007$, as does the
end-to-end contrast within the fixed geometry at $p=0.007$. Three of the four
intermediate single steps do not, which is why the argument rests on the
endpoints rather than on the step-by-step sequence.

Neither axis dominates the other and we do not claim one does. Over the range
tested the two effects are of similar size and opposite sign: $-0.027$ from
quadrupling the points inside a fixed geometry, $+0.024$ from going from twenty
units to thirty-two at $n=800$. The representative control declines slightly
with sample size as well, by $0.007$ with an interval containing zero, so the
fixed-geometry effect net of it is $-0.020$ with a 95 per cent interval of
$[-0.039, -0.000]$, which is real but marginal. The conclusion that survives is
weaker than a ranking of the two axes and more useful than one: a larger
calibration sample is not by itself a better one, and where the extra points go
decides whether the interval improves or degrades.

The reason is visible in the Gap column, which divides the error over the
domain by the out-of-fold error of the calibration sample itself. The
out-of-fold error moves little across the table, from 29.0 to 32.2 \Mgha{} over
the clustered cells and 33.6 at its widest once the representative control is
included, so the conformal quantile is an honest summary of the sample it was
built from in every cell. What the design changes is the error everywhere else.
The gap is 1.15 to 1.18 across the fixed
geometry and falls from 1.40 to 1.10 as units are added. The representative
control shows what its absence looks like: under simple random sampling the gap
is 1.01 at all three sizes and coverage is 0.915, 0.907 and 0.907, so there
sample size does what sample size is supposed to do, improving the map from
33.96 to 30.83 \Mgha{} and leaving the interval valid. Coverage falls as the
gap widens, with a correlation of $-0.98$ over the campaigns. The relation is
not an exact ordering: two pairs whose gaps differ by less than 0.05 sit the
other way round, so the gap explains the large differences and not the small
ones. The interval is not miscalibrated for its own calibration sample. It is
miscalibrated for the domain, and only a design that samples the domain can
reveal the difference.

Two qualifications. Primary units are drawn with replacement from one hundred
candidate clusters, so the realised number of distinct clusters is 7.7, 18.2 and
27.5 rather than 8, 20 and 32; the parenthesised figures in
Table~\ref{tab:sizes} are the realised expectations. The same mechanism means
the wider arm also densifies slightly as it widens, by about 12 per cent in
points per distinct cluster, which acts against the effect reported and makes
$+0.067$ a conservative bound rather than an inflated one.

No geometry reaches the nominal level. Widening the footprint improves coverage
without restoring it, which is the same conclusion Section~\ref{sec:intervals}
reaches from the theory. The spatial-block calibration, by contrast, returns
between 0.930 and 0.946 in every cell of the table, so the repair is insensitive
to both the size and the geometry that break the plain interval.

\begin{table}[!ht]
\centering
\caption{Sample-size and geometry sensitivity, 30 repetitions per cell, nominal
0.90, validation draw of 500 throughout. Units $\times$ pts gives the primary
units and the points drawn within each, with the realised expected number of
distinct units in parentheses, primary units being drawn with replacement. Gap
is the design-based RMSE divided by the out-of-fold RMSE of the calibration
sample, so it measures how much larger the error over the domain is than the
error the conformal quantile was built from. RMSE is in \Mgha{}; PICP, SD and
the confidence interval describe the CV+ interval, with leading zeros dropped on
the bounded columns; Block is the spatial-block calibration.
$^{a}$The 20 $\times$ 25 campaign is the design shared by both clustered arms
and is one set of 30 repetitions reported twice, so the arms meet at it rather
than crossing.}
\label{tab:sizes}
\footnotesize
\setlength{\tabcolsep}{2pt}
\input{tables/sizes}
\end{table}

\subsection{Where the improvement actually comes from}
\label{sec:ablation}

Two questions hide inside the phrase ``the geostatistical part''. Does the
variogram-derived representation help at all, and does tuning its parameters
help on top of that? Table~\ref{tab:ablation} separates them under the simple
random design, with everything else held fixed.

The first answer is yes. Removing the variogram-derived features, so that the
residual learner keeps the plain coordinates but loses the radial basis, the
kriging learner and the kernel smoother, costs $0.180$ \Mgha{}: the correction
still helps, at $-0.670$ against the uncorrected forest rather than $-0.849$,
but a fifth of the total gain is gone. Paired across the 30 repetitions the
variogram-derived information helps in 27 of them, with a 95 per cent interval
of $[+0.087, +0.273]$ that excludes zero. Table~\ref{tab:ablation} splits that
$0.180$ further: $0.172$ is the representation itself and $0.008$ is the
selection of its parameters.

The second answer is no. Selecting the family, practical range, binning,
estimator and anisotropy rather than fixing them at the conventional values is
worth $0.008$ \Mgha{}. The paired 95 per cent interval, $[-0.038, +0.023]$,
contains zero, and the selected recipe wins in only 8 of the 30 repetitions:
when it wins it wins by more than it loses by, which is how a mean this small
survives a minority of wins, but nothing here separates it from the conventional
values.

The two together are the honest description of what the geostatistics does
here. Encoding spatial structure through a variogram earns its place in the
model. Choosing the variogram's parameters from the data does not earn its
place in the accuracy, and we do not claim it does. What that selection buys is
stated in Section~\ref{sec:ablation}: it costs nothing on the criterion it is
run against, it changes the fitted model in a little under half the
repetitions, and it lets us report that the conventional recipe was adequate
for these data rather than assume it.

\begin{table}[!ht]
\centering
\caption{Decomposition of the improvement under the simple random design, 30
repetitions, $n=500$, RMSE in \Mgha{}. The third row removes the
variogram-derived features from the residual learner while keeping the plain
coordinates, isolating what the variogram representation contributes. The
fourth selects the variogram family, practical range, binning, estimator and
anisotropy out of fold. Wins count repetitions in which that configuration
beats the uncorrected forest.}
\label{tab:ablation}
\small
\begin{tabular}{lrrr}
\toprule
Configuration & RMSE & $\Delta$ vs RF & Wins \\
\midrule
Default random forest & 32.493 & n/a & n/a \\
Tuned base forest, no correction & 32.466 & n/a & n/a \\
Correction, no variogram-derived features & 31.824 & -0.670 & 27/30 \\
Correction, selected geostatistical hyperparameters & 31.644 & -0.849 & 28/30 \\
\bottomrule
\end{tabular}

\end{table}

The decomposition is in four parts and they sum exactly to the $0.849$ \Mgha{}
that separates the default forest from the deployed map. Two of the four are
real and two are not. Tuning the base forest is worth $0.027$, with a paired 95
per cent interval of $[-0.018, +0.073]$ that contains zero and wins in 18 of 30
repetitions, so it is not separable from doing nothing. The residual correction,
stripped of any variogram-derived information, is worth $0.642$,
$[+0.456, +0.828]$, in 25 of 30. The variogram-derived representation adds
$0.172$, $[+0.072, +0.272]$, in 26 of 30. Selecting the variogram's parameters
rather than fixing them adds $0.008$, $[-0.023, +0.038]$, in 8 of 30.

Three quarters of the gain is therefore the residual correction, a fifth is the
variogram-derived representation of space, and the two parameter-choosing steps,
the forest tuning and the geostatistical search, contribute nothing that a
paired interval can distinguish from zero.

We state the consequence plainly rather than leaving it to the reader. On this
dataset, selecting the geostatistical hyperparameters does not make the map more
accurate, and any claim that it does would not survive Table~\ref{tab:ablation}.
What the selection does deliver is different and, we think, still worth having.
It changes what the fitted model is, displacing the conventional recipe in 14,
12 and 11 of 30 repetitions across the three designs. The anisotropy axis is a
separate matter and we should not claim it as a selection result. Anisotropy is
adopted only on an out-of-fold improvement above 3 per cent, a parsimony
threshold applied inside the geometry estimator, and under clustered calibration
that threshold is never cleared: in all 30 repetitions the anisotropic candidate
collapses onto its isotropic twin and the two score identically, so the search
never faced a live choice. The refusal is the threshold's, not the search's, and
Section~\ref{sec:sims} shows we cannot attribute it to the design either. It is not selected against, because the
null correction of Eq.~(9) bounds the out-of-fold criterion below by the
uncorrected forest. That is a statement about the selection criterion, not a
guarantee of non-inferiority on an independent population. And
it converts a modelling assumption into a measured statement: we can now say
that the conventional variogram recipe was adequate for this dataset, which is
not something one can assert without having tested it.

\subsection{What the simulations show, including where they contradict us}
\label{sec:sims}

The claims above about the geometry estimator, and our reading of the clustered
result, need evidence rather than assertion. We therefore ran controlled
simulations on Gaussian fields with exponential covariance and known anisotropy.

Table~\ref{tab:sim1} gives recovery under simple random sampling with 500
points. On isotropic fields the estimator declares anisotropy in none of the
runs, so the parsimony rule of Section~\ref{sec:geohyper} holds, and whenever it
does declare anisotropy the recovered major axis is exactly correct. Its
weakness is power rather than specificity: a ratio of $0.20$ is detected in 92
to 100 per cent of runs, $0.35$ in 29 to 58 per cent, and $0.60$ never. The
errors under this design are therefore false negatives, and the field results
should be read as a lower bound on how often directional structure is present.
Specificity is not perfect everywhere: under clustered sampling the estimator
declares anisotropy on a small fraction of isotropic fields
(Table~\ref{tab:sim2}), so the guarantee is one-sided only for the
representative case.

\begin{table}[!ht]
\centering
\caption{Simulation S1. Proportion of runs in which anisotropy is declared, by
true ratio and signal-to-noise ratio, under simple random sampling with 500
points. A true ratio of 1 is an isotropic field.}
\label{tab:sim1}
\small
\begin{tabular}{rrr}
\toprule
True ratio $\varrho$ & Signal-to-noise & Anisotropy declared \\
\midrule
1.00 & 2 & 0.00 \\
0.60 & 2 & 0.00 \\
0.35 & 2 & 0.29 \\
0.20 & 2 & 0.92 \\
1.00 & 8 & 0.00 \\
0.60 & 8 & 0.00 \\
0.35 & 8 & 0.58 \\
0.20 & 8 & 1.00 \\
\bottomrule
\end{tabular}

\end{table}

Table~\ref{tab:sim2} tests our own explanation of the clustered result and does
not support it. Holding the true anisotropy fixed and varying only the sampling
design, anisotropy is declared in 58.3 per cent of simple random runs and 50.0
per cent of clustered runs, a difference that is not distinguishable
(Fisher exact test, $p = 1$). Clustering therefore does not prevent the
estimator from recovering direction at this sample size, and the refusal
observed in the field experiment cannot be attributed to the design in the way
we first suggested. The likelier readings are that the residual field of this
dataset is close to isotropic at this scale, or that an anisotropic geometry
simply did not lower out-of-fold risk enough to win the profile search. We
report this because it is the kind of post hoc explanation that is easy to
assert and, as it turns out, wrong.

\begin{table}[!ht]
\centering
\caption{Simulation S2. Proportion of runs declaring anisotropy when the true
geometry is held fixed and only the sampling design changes.}
\label{tab:sim2}
\small
\begin{tabular}{lrr}
\toprule
Sampling design & True ratio $\varrho$ & Anisotropy declared \\
\midrule
Clustered & 0.35 & 0.500 \\
Simple random & 0.35 & 0.583 \\
Clustered & 1.00 & 0.083 \\
Simple random & 1.00 & 0.000 \\
\bottomrule
\end{tabular}

\end{table}

Finally, we checked the interval construction itself, and this is where the
simulations earn their place. One field realisation is generated over the whole
domain; the target points are always uniform over it and only the calibration
design changes. The field is stationary by construction, with one fixed anisotropic geometry
shared by both arms, so anything the contrast produces is attributable to the
geometry of the calibration sample and to nothing else.

Under simple random calibration CV+ attains empirical coverage of $0.912$
against a nominal $0.90$, comfortably above the $1-2\alpha$ bound of $0.80$,
which for $K=5$ and $n=300$ carries a further deduction of about $0.026$ under
Theorem 4 of \citet{barber2021}, and mildly conservative, as the theory
predicts. Under clustered calibration it falls to $0.811$.

The clustered figure is the informative one. A stationary field is enough: no
non-stationarity, no covariate shift beyond what the sampling geometry itself
induces, and coverage still fails. The mechanism is spatial extrapolation. Calibration
residuals are collected where the model interpolates between nearby points,
while the target population includes locations far from every cluster, where it
must extrapolate and where its errors are larger than any the conformity scores
saw. The simulation undercovers more than the biomass surface does, at $0.811$
against $0.861$, which is consistent with the simulated clusters occupying a
smaller share of their domain than the field sample does of the Amazon
rectangle.

\section{Discussion}

The uncertainty result is the one we would defend first. Extending design-based
assessment from error to coverage is conceptually slight, three lines in
Section~\ref{sec:ht}, but it changes what can be checked, and the check has
teeth. Intervals that look entirely reasonable, and that are demonstrably
well behaved and sharp under representative sampling, cover only 0.861 of the
population once the calibration sample is clustered. That is the same regime in
which \citet{wadoux2021} found point-accuracy estimators to fail, and it fails
for a related reason: the sample is not representative of the domain the map has
to cover. The practical recommendation follows directly. Assess coverage
design-based, and when clustering is suspected use a spatial-block calibration,
accepting wider but valid intervals.

We are careful about how strong a claim this supports. CV+ carries a
finite-sample marginal guarantee of $1-2\alpha$ less a term of order $K/n$ under
exchangeability
\citep{barber2021}, and our simulations confirm that it is attained. Clustered
sampling violates exchangeability, so no guarantee survives it, and what
Section~\ref{sec:intervals-results} reports for the spatial-block variant is
restored \emph{empirical} coverage in this experiment rather than a theorem. The
comparison with quantile regression forests is likewise qualified: the two
methods sit at different effective coverage levels, 0.907 against 0.937, and the
quantile forest holds 0.928 under clustered calibration where CV+ falls to
0.861, so the
16 per cent width advantage should be read alongside that difference rather than
as a like-for-like sharpening.

The point predictor is the more modest contribution, and the ablation is what
disciplines the claim. Table~\ref{tab:ablation} attributes three quarters of
the gain to the residual correction and a fifth to the variogram-derived
representation, leaving the two parameter-choosing steps with nothing a paired
interval separates from zero. We therefore do not argue that moving the variogram into the
hyperparameter space makes maps better. We argue something narrower: that it can
be done without penalty on the selection criterion, since Eq.~(9) bounds it
below by the uncorrected
forest; that it changes the fitted model, displacing the conventional recipe in
14, 12 and 11 of 30 repetitions; that under clustered calibration no anisotropic
geometry ever clears the parsimony threshold that would let the search consider
one; and that it replaces an assumption with a
measurement, letting us report that the conventional recipe was adequate here
instead of taking it on faith.

Two further boundaries deserve stating, because overstating them would repeat
the error \citet{wadoux2021} criticised. The corrector the data selects is
usually a residual forest, and ordinary kriging wins outright in three of 90
repetitions, so the geostatistics contributes the representation and the
regularisation rather than the final estimator. And our own explanation of the
clustered anisotropy result did not survive testing: Section~\ref{sec:sims}
holds the true geometry fixed and varies only the design, finding no detectable
difference in recovery, so the refusal cannot be attributed to the design in the
way we first supposed.

\paragraph{What the correction does not fix} The point gain is incremental, of order 0.85 to 1.10
\Mgha{}. The anisotropy model is geometric and global, one angle and ratio for
the whole domain, and the geometry estimator has limited power against mild
anisotropy. The covariate-shift calibration does not help at all, for the
structural reason that reweighting cannot manufacture residuals a clustered
sample never observed. The reference surface is a remote-sensing product treated
as a finite population, so conclusions concern reconstructing that benchmark
rather than agreement with field measurements. And all field results come from a
single dataset, although Sections~\ref{sec:distance} and \ref{sec:sizes} vary
the sampling design and the calibration size within it.

\section{Limitations}
\label{sec:limits}

Four limits bound what the results above support.

The evidence comes from one surface, and it is not a soil survey. The
benchmark is above-ground biomass, so the paper is framed as spatial
predictive mapping rather than digital soil mapping. The derivation applies to
any map of a spatial variable built from a point sample, and soil properties
are an obvious case, but a pedological campaign brings its own sampling
traditions and covariate structure and we have not tested one. The benchmark of \citet{baccini2012} is a
remote-sensing product with its own error, treated here as a finite population
because \citet{wadoux2021} treated it that way and comparability matters more
than realism for this purpose. Nothing in the experiment speaks to how the
undercoverage would scale on a different variable, region or resolution.

Sample size is varied for the coverage result and not for the rest.
Section~\ref{sec:sizes} spans calibration sizes from 200 to 800 on two design
axes, so a small-sample explanation for the clustered shortfall is excluded over
that range, and the spatial-block calibration holds its coverage in every cell.
Its width does not: the repaired interval narrows from 144.9 to 136.0 \Mgha{}
as points are added to a fixed geometry and from 167.3 to 129.1 as units are
added, so what the sweep establishes about the repair is that it stays
conservative, not that its cost is stable. The point-accuracy comparison, the
ablation, the estimand comparison and the counts of how often the search
displaces the conventional recipe or declines an anisotropic geometry all still
come from $n = 500$ alone.

The two repairs are compared empirically, not like for like. The CV+ row is a
different construction from the covariate-shift and spatial-block rows, which
modify a pooled calibration, and only CV+ carries a finite-sample bound. Their
coverage figures are therefore descriptive.

The spatial-block result is a repair that worked here. It reached $0.946$
against a nominal $0.90$ with intervals about half as wide again, which is
conservative rather than calibrated. It held between $0.930$ and $0.946$ across
every design, size and geometry in Table~\ref{tab:sizes}, which is more than a
single point but is still one dataset, and we have no theorem that says it will
recover coverage on another.

\section{Conclusion}

Coverage is a population parameter of the same kind as the mean squared error,
and it inherits the same design-unbiased estimator. Saying so requires no new
theory and
reveals something: cross-fitted conformal intervals that are calibrated and
sharp under representative sampling under-cover badly once the calibration
sample is clustered, a failure no exchangeability-assuming check can detect, and
a spatial-block calibration returns coverage above nominal in this experiment.
Enlarging the calibration sample does not help, and spent inside a fixed set of
primary units it hurts while the map improves, because what the interval lacks
is not data but reach: the error grows with distance from the calibration sample and the
interval width does not follow it. Map uncertainty, like map accuracy, should be
confronted with the sampling design, and a campaign that wants believable
intervals should buy locations rather than points.

The geostatistical decisions behind the point map can meanwhile be treated as
hyperparameters and selected out of fold rather than fixed by habit. The
correction as a whole lowers design-based RMSE under every design, but choosing
the variogram's parameters rather than fixing them by convention contributes
nothing to that on these data. It costs
nothing, it changes what the model does, it declines parameters the data do not
reward, and it lets a modeller state that the conventional recipe was adequate
rather than assume it. The corrector the model finally prefers is a
forest, but it is a forest that sees the landscape through a variogram.

\appendix
\section{Full computational specification}
\label{app:spec}

Every setting needed to reproduce the experiments is stated here; all values are
the defaults in the released code unless noted.

\paragraph{Cross-fitting and selection} Two fold partitions are used and they
are not the same. The selection folds are $K=5$ assigned by
\texttt{caret::createFolds}, which stratifies on the response; fold membership
there depends on $y_i$, which is acceptable because those folds only rank
candidate configurations. The calibration partition used by the nested CV+ of
Section~\ref{sec:intervals} is a plain random permutation into $K=5$ blocks,
drawn independently of the response, because the exchangeability the bound needs
would not survive response-stratified assignment. The shrinkage grid is
$\lambda \in [-0.25, 0.60]$ in steps of $0.05$, 18 values, for $n \ge 300$;
narrower grids apply at smaller $n$. Selection minimises Eq.~(8) over the
correction library crossed with that grid, with $\lambda=0$ included.

\paragraph{Base forest profiles} Six fixed profiles are compared on identical
folds: \file{default_rf}, \file{more_trees}, \file{deeper_trees},
\file{subsample}, \file{wide_mtry}, \file{wide_deep_subsample}. Across
them the number of trees takes values in $\{500,1000,1500,2000\}$,
\texttt{mtry} in $\{5,9,14\}$ for $p=28$, minimum node size in $\{1,5\}$ and
sample fraction in $\{0.632,0.8,1\}$; residual forests use 300 or 500 trees.
Where explicit grid tuning is enabled instead, the base grid is capped at 18
candidates and the residual grid at 12, scored out of bag.

\paragraph{Correction library} Residual random forest; residual Cubist with 10
committees and 0 neighbours; a single-hidden-layer network; ordinary kriging of
residuals; a radial basis or kernel smoother; the unweighted mean of the
residual learners; and bias, linear, quadratic and stacked recalibrations. The
radial basis places centres by $k$-means on the coordinates, their number
derived from the domain area, the practical range and the nugget ratio and
clamped to $[8,48]$, with bandwidth $b=m_b\rho$.

\paragraph{Variogram} Empirical variogram with the robust estimator of
\citet{cressie1980}, cutoff at half the maximum pairwise distance, 15 lags.
Exponential, spherical, Gaussian and Mat\'ern models are each fitted by weighted
least squares to the same empirical variogram, and the simplest family within a
relative tolerance of $0.05$ of the best weighted residual sum of squares is
retained. The practical range is the distance at which the fitted model first
reaches 95 per cent of its sill, evaluated numerically so that it means the same
thing for every family.

\paragraph{Search axes} Three variogram configurations join the default: cutoff
$0.35$ with 15 lags, cutoff $0.65$ with 20 lags, and the classical
\citet{matheron1963} estimator at the default binning. The range multiplier
takes $m_b \in \{0.25,0.5,1.0\}$, with $0.5$ the default arm. The anisotropic
variant is offered for the two leading profiles; candidate geometries are angles
$\{0,22.5,\dots,157.5\}$ degrees crossed with ratios
$\{1,0.75,0.55,0.40,0.25\}$, scored by 5-fold out-of-fold error of a kernel
smoother whose bandwidth is $0.15$ times the median pairwise distance in the
transformed space, and accepted only on an improvement exceeding 3 per cent. The
threshold is a parsimony rule rather than a test: the anisotropic geometry costs
two extra parameters, and a margin of this size keeps it out unless the
out-of-fold gain is larger than the run-to-run variation we observed in the
isotropic fits.

\paragraph{Gate and intervals} The gate length multiplier is chosen by spatially
blocked cross-validation with four $k$-means blocks, the candidate set including
the value that disables the gate. Intervals are formed with $K=5$ and $\alpha=0.10$ over the calibration partition
described above, which is drawn independently of the selection folds and of the
response; the selection is re-run inside each of those folds, as
Section~\ref{sec:intervals} sets out. The covariate-shift variant estimates density ratios with a forest
classifier separating calibration points from a sample of the population
covariates, clipping weights at 20; the spatial-block variant refits on five
$k$-means blocks. The quantile regression forest baseline uses \texttt{ranger}
with \texttt{quantreg = TRUE} and the same predictors.

\paragraph{Sampling designs} Simple random sampling draws $n=500$ cells without
replacement; the systematic design draws a regular grid; the clustered design is
two-stage with 20 primary units chosen over 100 $k$-means centres and 25 points
each. Primary units are drawn with replacement, so the realised number of
distinct units is smaller than the nominal one, 18.2 in expectation for 20 and
7.7 and 27.5 for the 8 and 32 used in Section~\ref{sec:sizes}. The validation
sample is an independent simple random sample of $n=500$ drawn from the same
frame without excluding calibration cells, an overlap of about one cell in
2600 at that size. The Population estimand scores a fixed 20{,}000-cell
reference subset drawn once with its own seed, also without exclusion. Spatial
block cross-validation uses complete-linkage clustering of coordinates cut at
350 km.

\paragraph{The sensitivity campaigns} Section~\ref{sec:sizes} varies the
calibration size over 200, 500 and 800 with 30 repetitions per cell and the
validation draw fixed at 500, which the runner takes from a separate setting so
that the precision of the measurement does not move with the quantity measured.
Under the random design the size is the number of cells drawn. Under the
clustered design it is the product of the primary-unit count and the points
within each, and the two are varied separately: 20 units by 10, 25 and 40
points, and 8, 20 and 32 units by 25 points. The 20 by 25 cell is shared by both
arms and is reported once. Section~\ref{sec:distance} re-runs three of these
designs with 10 repetitions each and an additional per-point record: the
coordinates of every validation point, its distance to the nearest calibration
point, the observation, the prediction and both interval bounds. That record is
written only when explicitly requested, since it is three orders of magnitude
larger than the aggregated outputs.

\paragraph{Preprocessing, missing data, seeds} Covariates and response enter
untransformed. Rows with non-finite response or coordinates are dropped before
variogram fitting, and constant or duplicated predictor columns are removed
before fitting. Every \texttt{ranger} call receives an explicit seed and
\texttt{ranger} 0.18 and later is reproducible across thread counts; the
estimand campaign draws repetition $i$ with seed $123 + 1000(i-1)$, while the
30-repetition campaigns seed the stream once with $123$ and draw the
repetitions in sequence. The internal
random state is restored on exit with \texttt{withr}, so calling the estimator
does not alter the user's stream.

\paragraph{Software and scripts} R 4.6.0 with \texttt{ranger}
\citep{wright2017}, \texttt{gstat} \citep{pebesma2004}, \texttt{caret},
\texttt{Cubist} and \texttt{nnet}. The estimator is the \texttt{GeoVersa}
package. Every table and figure in this paper is produced by
\file{run_dbtarf_wadoux_comparison.R} (accuracy campaigns),
\file{run_dbtarf_protocols_comparison.R} (estimands),
\file{simulations.R} (Section~\ref{sec:sims}), \file{tables_manuscript.R}
(all tables), \file{maps_compute.R} with \file{maps_plot.R} (the maps),
\file{figures_schematic.R}, \file{figures_protocols.R},
\file{figures_uncertainty.R} and \file{figures_distance.R} (the remaining
figures). The campaign sets are driven by \file{run_all_campaigns.sh} for the
headline designs, \file{run_size_sweep.sh} for Section~\ref{sec:sizes} and
\file{run_point_dump.sh} for Section~\ref{sec:distance}; \file{freeze_check.sh}
is the pre-flight that gates a launch, and refuses to run against a working tree
with uncommitted changes. Every campaign directory records the git commit,
whether the tree was dirty, the R version, a timestamp and the sampling
parameters that governed it, and the table generator drops any campaign whose
recorded calibration size is not the size the table claims for it.


\section*{Data and code availability}

No new primary data were collected for this study.

\paragraph{Software} The estimator, the campaign runners and every script that
produces a table or figure in this paper are released as the \texttt{GeoVersa}
package for R at \url{https://github.com/HugoMachadoRodrigues/GeoVersa} under
the MIT licence. The results reported here were produced at commit
\texttt{5db0ec3}. An archived release carrying a persistent identifier is
deposited at Zenodo, \textsc{doi} to be inserted at acceptance.

\paragraph{Benchmark materials} The analyses read the public materials of
\citet{wadoux2021}, distributed by those authors at
\url{https://github.com/AlexandreWadoux/SpatialValidation}. These comprise the
Amazonian above-ground live woody biomass surface of \citet{baccini2012} at 1 km
resolution, 28 covariate rasters and the study polygon. They are the
intellectual property of their respective authors and publishers and are
\emph{not} redistributed with this work. Readers obtain them directly from the
source repository, under whatever terms those authors set, and
\file{data/README.md} in our repository gives the command and the directory the
code expects. Our repository excludes them by rule rather than by omission: the
version-control configuration excludes the \file{data}, \file{external} and
\file{results} directories in their entirety and re-admits only their README
files, so no third-party file can be committed by oversight. A report
enumerating every file the repository distributes, together with the search
used to establish it, is generated by \file{tools/data-provenance-report.sh} and
can be re-run by any reader from a clone.

\paragraph{Derived results} The campaign outputs behind every table and figure,
comprising thirteen campaigns of thirty repetitions and three of ten, amount to
roughly 833 MB and are not distributed. They are reproducible from the sources
above: \file{code/run_all_campaigns.sh} produces the three headline designs,
\file{code/run_size_sweep.sh} the size and geometry sweep of
Section~\ref{sec:sizes}, and \file{code/run_point_dump.sh} the per-point records
behind Fig.~\ref{fig:distance}. Every results directory records the git commit,
whether the working tree was clean, the R version, a timestamp and the sampling
parameters that governed the run, and the table generator refuses any campaign
whose configuration lacks a commit. Complete computational settings, including
seeds, are given in Appendix~A.

\section*{Declaration of generative AI and AI-assisted technologies in the
manuscript preparation process}

During the preparation of this work the authors used Grammarly for spelling and
grammar checking, and an AI assistant (Claude, Anthropic) to draft manuscript
text, implement analysis code and produce figures, under the authors' direction.
All study design, methodological decisions, interpretation of results and
verification of the analysis are the authors' own. After using these tools the
authors reviewed and edited the content as needed and take full responsibility
for the content of the publication.

\section*{CRediT authorship contribution statement}

\textbf{Hugo Rodrigues}: Conceptualization, Methodology, Software, Formal
analysis, Investigation, Data curation, Visualization, Writing -- original
draft. \textbf{Golmar Golmohammadi}: Supervision, Validation, Writing -- review
and editing.

\section*{Funding}

This research did not receive any specific grant from funding agencies in the
public, commercial, or not-for-profit sectors.

\section*{Declaration of competing interest}

The authors declare no known competing financial interests or personal
relationships that could have appeared to influence the work reported in this
paper.

\bibliographystyle{elsarticle-harv}
\bibliography{references}

\end{document}
